\chapter{ESPECIFICAÇÃO DO PROTÓTIPO}
\label{cap:prototipo}

Este capítulo apresenta a especificação conceitual do protótipo elaborada a partir dos resultados da revisão sistemática e da fundamentação pedagógica. O artefato produzido consolidou princípios de projeto, requisitos funcionais e não funcionais, critérios para fontes de dados e modelos, um protocolo de avaliação e uma arquitetura de referência.

A especificação não foi tratada como evidência de uma aplicação funcional. Seu resultado consistiu em transformar as evidências da literatura em decisões técnicas e pedagógicas documentadas, auditáveis e coerentes com o escopo do trabalho.

\section{Princípios de Projeto}

A especificação estabelece que o protótipo deverá apoiar a decisão pedagógica sem substituir o professor. Essa decisão de projeto implica distinguir, na representação e na interpretação dos resultados, o desempenho observado em uma atividade, a proficiência estimada por um modelo, a competência mobilizada em um contexto mais amplo e a aprendizagem como possível mudança ao longo do tempo. Esses conceitos não são intercambiáveis, e a passagem de um para outro não poderá ser presumida apenas pela produção de um indicador computacional.

As estimativas deverão ser acompanhadas de explicações, limitações e condições de validade compreensíveis para seus destinatários. A revisão humana deverá permanecer integrada ao uso das recomendações, de modo que o professor possa confrontar as evidências com o contexto da turma e decidir se há fundamento para uma intervenção. Nesse sentido, a saída do protótipo será tratada como apoio à análise, e não como classificação definitiva ou prescrição pedagógica.

A proteção dos dados educacionais constitui outro princípio da especificação. O uso deverá estar associado a uma finalidade explícita, com minimização dos dados pessoais e controle de acesso compatível com essa finalidade. Também deverá ser possível reproduzir e auditar os procedimentos de preparação, treinamento e avaliação por meio do registro das versões, dos parâmetros e das dependências relevantes. Por fim, a especificação rejeita rótulos permanentes ou deterministas sobre os estudantes, pois uma estimativa situada não deve ser convertida em atributo fixo da pessoa.

\section{Requisitos Funcionais}

Foram definidos os seguintes requisitos funcionais para o protótipo:

\begin{enumerate}
    \item importar dados documentados e relacionados a objetivos curriculares;
    \item organizar evidências de desempenho por habilidade ou descritor;
    \item produzir estimativas de proficiência acompanhadas de incertezas e limitações;
    \item apresentar resultados individuais e agregados ao professor;
    \item explicar as variáveis relacionadas a cada estimativa;
    \item permitir revisão humana das recomendações produzidas;
    \item exportar dados, parâmetros e resultados para auditoria.
\end{enumerate}

A especificação também estabeleceu que indicadores computacionais não representam, isoladamente, diagnósticos de aprendizagem. Sua interpretação depende da natureza da atividade, da qualidade dos dados e do contexto pedagógico.

\section{Requisitos Não Funcionais}

Os requisitos não funcionais especificam condições de qualidade e de governança que deverão ser satisfeitas por uma futura implementação. Eles não constituem evidência de que tais propriedades já tenham sido alcançadas. A reprodutibilidade deverá ser assegurada pelo registro das versões dos dados, dos parâmetros, dos modelos e das dependências. Complementarmente, a rastreabilidade deverá permitir associar cada resultado à fonte de dados utilizada e à versão do experimento que o produziu. Essa associação é necessária para que uma conclusão possa ser examinada sem perder sua procedência.

A testabilidade deverá ser favorecida pela separação entre ingestão, preparação, modelagem, avaliação e apresentação. A portabilidade, por sua vez, exigirá um ambiente de execução documentado e passível de reprodução. Essas duas condições dizem respeito à possibilidade de verificar e repetir os procedimentos, e não à escolha de uma tecnologia ou de uma interface específica.

Os requisitos também abrangem a relação do protótipo com seus usuários e com os dados tratados. A apresentação deverá ser acessível, compreensível e compatível com diferentes perfis de usuário. A privacidade deverá orientar a minimização de dados pessoais e o controle de finalidade e de acesso. Por fim, a transparência deverá exigir que as limitações dos modelos e as condições de validade das estimativas sejam comunicadas junto aos resultados, evitando que a forma de apresentação atribua uma certeza que o procedimento não sustenta.

\section{Critérios para Fontes de Dados}

A especificação toma como referência a possibilidade de trabalhar com microdados de avaliações educacionais, como SAEB e PISA, e com dados sintéticos destinados à verificação do \textit{pipeline}. Essa referência não representa a seleção de uma base definitiva. Trata-se de uma decisão metodológica preliminar, que mantém a escolha da fonte condicionada à documentação disponível, ao problema de pesquisa e às condições de execução que vierem a ser definidas.

A seleção de uma fonte deverá considerar, inicialmente, sua licença e suas condições de uso, bem como a existência de dicionário de dados e de documentação metodológica. Também deverá examinar a unidade de análise e sua relação com o problema de pesquisa, a qualidade dos metadados e o tratamento previsto para valores ausentes. A representatividade da população deverá ser analisada em relação ao contexto para o qual se pretende interpretar os resultados, sem ser presumida a partir da simples disponibilidade da base. Do mesmo modo, a presença de variáveis relacionadas ao desempenho e ao contexto educacional será relevante, mas não garantirá, por si só, que seja possível associá-las a competências ou descritores curriculares. Essa associação deverá ser demonstrada e delimitada pela documentação e pelo desenho da avaliação.

Por fim, a análise da fonte deverá incluir os riscos de identificação, discriminação e uso indevido. Esses riscos são condições de governança da execução, e não apenas atributos técnicos do arquivo de dados. Dados sintéticos poderão ser usados para testar fluxos de processamento, interfaces e regras de validação. Entretanto, não constituirão evidência de desempenho educacional, de validade preditiva ou de eficácia pedagógica, nem substituirão uma base empírica quando a investigação exigir inferências sobre estudantes ou contextos reais.

\section{Alternativas de Modelagem}

A revisão identificou uso recorrente de \textbf{\textit{Random Forest}}, máquinas de vetores de suporte (SVM), árvores de decisão, regressão logística e redes neurais. Na especificação, \textit{Random Forest} e SVM também foram ancoradas em suas fontes metodológicas primárias \cite{Breiman2001,CortesVapnik1995}. Modelos supervisionados simples e interpretáveis foram mantidos como referências iniciais de comparação, condicionando o uso de modelos mais complexos à demonstração de ganho e à possibilidade de interpretação adequada.

\textbf{\textit{Random Forest}} e SVM foram registrados como alternativas de modelagem, não como algoritmos definitivamente implementados. A escolha entre eles depende da base selecionada, da distribuição das classes, do tipo de variável, do custo computacional e dos requisitos de explicabilidade.

Em uma implementação com scikit-learn, identificadores longos, como \texttt{RandomForest\allowbreak Classifier}, devem ser documentados sem comprometer a legibilidade do texto ou a reprodução do código.

\section{Protocolo de Avaliação}

O protocolo de avaliação constitui uma condição para a execução futura da modelagem e para a interpretação dos seus resultados. Ele deverá separar o que será observado no comportamento técnico dos modelos daquilo que exigiria evidência educacional adicional. Essa separação começa pela definição explícita dos conjuntos de treino, validação e teste e pela prevenção de vazamento de informação entre preparação, modelagem e avaliação. Os modelos candidatos deverão ser comparados com referências simples, para que qualquer ganho seja interpretado em relação a uma base de comparação, e as métricas deverão ser analisadas por classe, além das medidas agregadas, quando a tarefa permitir essa decomposição.

Como descrição complementar dos erros e acertos, o protocolo prevê o uso de matriz de confusão, precisão, revocação e medida F1. A área sob a curva ROC e a calibração somente deverão ser calculadas quando forem compatíveis com a tarefa e com o tipo de saída produzido \cite{Fawcett2006,NiculescuMizilCaruana2005}. Também deverão ser examinadas a estabilidade das explicações e possíveis diferenças entre grupos, sem converter essas análises em conclusões causais não sustentadas pelo desenho do estudo. Por fim, a execução deverá registrar as limitações de generalização relacionadas à fonte de dados, à população observada e às condições de aplicação.

O protocolo distingue desempenho observado, proficiência estimada, competência e aprendizagem. O desempenho observado corresponde aos registros produzidos nas atividades ou avaliações consideradas. A proficiência estimada é uma saída inferida a partir desses registros e das variáveis admitidas pelo modelo. Competência e aprendizagem, contudo, não serão tratadas como equivalentes automáticos dessas duas grandezas: sua análise depende do contexto educacional, do construto adotado e, no caso da aprendizagem, de evidência de mudança ao longo do tempo.

Essa distinção também separa a avaliação técnica da eficácia pedagógica. As métricas preditivas descrevem o comportamento do modelo em uma base e sob determinadas condições de avaliação. Elas não demonstram, por si só, melhoria da aprendizagem, adequação de uma intervenção escolar ou alinhamento curricular validado. Tais conclusões exigiriam um desenho empírico próprio, com condições de comparação e observação educacional compatíveis com a pergunta investigada.

\section{Explicabilidade e Participação Docente}

SHAP e LIME foram incluídos como alternativas para explicar modelos, com referência às formulações originais de Lundberg e Lee e de Ribeiro, Singh e Guestrin \cite{LundbergLee2017,RibeiroSinghGuestrin2016}. A especificação condicionou seu uso à estabilidade das explicações, ao custo computacional e à possibilidade de tradução dos resultados para uma linguagem curricular compreensível.

O professor permaneceu como responsável por interpretar as evidências, relacioná-las ao contexto da turma e decidir intervenções. O protótipo foi especificado como instrumento de apoio, e não como mecanismo autônomo de classificação ou prescrição pedagógica.

\section{Arquitetura de Referência}

\textbf{Ingestão} corresponde à leitura das fontes, à validação inicial de seus formatos e à preservação das informações de procedência. Em seguida, a etapa de \textbf{preparação} deverá tratar valores ausentes, codificar e transformar variáveis e documentar as decisões tomadas sobre esses dados. A \textbf{modelagem} abrangerá o treinamento e a comparação dos modelos candidatos de acordo com a base selecionada e com o protocolo de avaliação.

A etapa de \textbf{avaliação e explicabilidade} deverá calcular as métricas, analisar os erros e gerar explicações cuja estabilidade e limitações serão examinadas. Por fim, a \textbf{apresentação} deverá disponibilizar indicadores, incertezas e limitações de modo compreensível para a análise docente. Essa arquitetura de referência organiza responsabilidades e pontos de verificação; não descreve um sistema funcional já disponível, nem fixa a tecnologia ou a forma de implantação.

A separação dos componentes deverá favorecer a testabilidade, a rastreabilidade e a substituição controlada de técnicas. Ela também permite localizar em qual etapa uma decisão, uma transformação ou uma limitação foi introduzida. A arquitetura não impõe uma interface específica e poderá ser materializada como \textit{notebook}, aplicação local, API ou sistema web. Qualquer dessas formas de execução deverá, entretanto, preservar os requisitos de auditoria, interpretação, privacidade e participação docente estabelecidos neste capítulo.

\section{Resultado da Especificação}

O resultado deste capítulo foi um artefato de especificação que conectou as evidências da revisão às decisões de projeto. Foram definidos o propósito do protótipo, seus limites, os requisitos mínimos, os critérios de seleção de dados e modelos, o protocolo de avaliação e a arquitetura de referência.

Essa entrega permitiu concluir o objetivo de transformar a literatura analisada em uma base técnica e pedagógica estruturada, sem confundir uma especificação validada documentalmente com uma implementação funcional ou com evidência de eficácia em sala de aula.

%--------- FIM ESPECIFICAÇÃO DO PROTÓTIPO ------------
